\documentclass[11pt]{beamer}

%\usetheme{default}
%\usecolortheme{default}
\beamertemplatenavigationsymbolsempty
\setbeamertemplate{footline}[frame number]

\usepackage[T1]{fontenc}
\usepackage[utf8]{inputenc}
\usepackage{listings}
\usepackage{graphics}
\usepackage{color}
\usepackage{xcolor}
\usepackage{hyperref}

\hypersetup{
    colorlinks = true
}

\AtBeginSection[]{
\begin{frame}{Contents}
    \tableofcontents[currentsection]
\end{frame}
}

\lstdefinestyle{shell}{
  basicstyle=\ttfamily\footnotesize,
%  identifierstyle=\ttfamily
  breaklines=true,
  showstringspaces=false,
  frame=single,
  framerule=0.4pt,
  xleftmargin=0.5em,
  xrightmargin=0.5em,
  columns=flexible,                   % Removes fixed-column spacing
  showstringspaces=false                  % Hides space markers in strings
}
\lstset{style=shell}

\title{Git and GitHub for Mathematicians}
\subtitle{Lecture 1 --- The Shell and Git Basics}
\author{}
\date{}

\newcommand{\exercise}[1]{
  \begin{frame}{Exercise}
    #1
  \end{frame}
}

\begin{document}

\begin{frame}
  \titlepage
\end{frame}

% ------------------------------------------------------------------

\begin{frame}{Why version control?}
  Typical pain points in writing a paper:
  \begin{itemize}
    \item \texttt{paper\_v2\_final\_REALLY\_FINAL.tex}
    \item Changes overwritten by a coauthor
%    \item ``It worked yesterday, what did I change?''
    \item Sending \texttt{.zip} files by email
    \item No record of \emph{why} a given change was made
  \end{itemize}
  \vfill
  Version control solves all of this.
\end{frame}

\begin{frame}{Git vs. GitHub}
  \textbf{Git}
  \begin{itemize}
    \item A \emph{distributed} version control system
    \item Runs locally on your computer
    \item Free software, works offline
  \end{itemize}
  \vfill
  \textbf{GitHub}
  \begin{itemize}
    \item A commercial website that hosts Git repositories
    \item Adds collaboration features (issues, pull requests, \ldots)
    \item Alternatives: GitLab, Bitbucket, Codeberg, \ldots

      Univie offers git hosting, but their service is pretty primitive
  \end{itemize}
  \vfill
  \alert{They are not the same thing.}
  \vfill
  Please, log out from GitHub now
\end{frame}

\begin{frame}{GUIs for git}
  Git is a textual program, i.e., it works the
  \emph{shell} (text console)

  \vfill
  A few graphical programs built over git are available
  \begin{itemize}
  \item gitk  (available on Univie computers)
  \item sourcegit
  \item integration in various editors (VSCode, Emacs, Vim\ldots)
  \item and others \url{https://git-scm.com/tools/guis}
  \end{itemize}

  We do not cover the usage of such tools, but they might be useful,
  especially for the content of the second lecture.
\end{frame}

\begin{frame}{Bibliography}
  \begin{itemize}
    \item \emph{Pro Git} (beyond our purpose) \url{https://git-scm.com/book}
    \item Git cheat sheet \url{https://git-scm.com/cheat-sheet}
    \item Oh My Git \url{https://ohmygit.org}
    \item Git Immersion \url{https://gitimmersion.com}
    \item Learn Git Branching \url{https://learngitbranching.js.org}
    \end{itemize}

    \vfill

    But mostly\ldots learn by doing!
\end{frame}

\begin{frame}{Contents}
  \tableofcontents
\end{frame}




% ------------------------------------------------------------------
\section{A minimal shell survival kit}

\begin{frame}{The shell}
  A \emph{shell} is a program that reads commands and runs them.
  \begin{itemize}
    \item The \emph{prompt} is where you type
    \item At all times there is a \emph{current working directory}
    \item Commands usually act relative to that directory
  \end{itemize}
  \vfill
  You only need a handful of commands for this course.
\end{frame}

\begin{frame}[fragile]{Navigating the filesystem}
\begin{lstlisting}
ls                     # list files
ls -l                  # long listing
cd some/dir            # change directory
cd ..                  # go up one level
cd                     # go to the home directory
mkdir directory        # create a directory
\end{lstlisting}
%ls -a             # include hidden files
%cd -         # go back to previous directory
\vfill
Content after \texttt{\#} are \emph{comments} and shouldn't be typed in
the shell
  \vfill
  Useful tricks:
  \begin{itemize}
    \item \texttt{Tab} completes
    \item Up/down arrows browse the command history
%    \item \texttt{man cmd} or \texttt{cmd -{}-help} for documentation
  \end{itemize}
\end{frame}

\exercise{
  \begin{itemize}
  \item Create a directory for the exercises of this course
  \item Familiarize with this new commands
  \end{itemize}
}

% ------------------------------------------------------------------
\section{Configuring Git}

\begin{frame}[fragile]{First-time setup}
  Tell Git who you are (only once per machine):
\begin{lstlisting}
git config --global user.name "Your Name"
git config --global user.email "you@example.com"
\end{lstlisting}
  \vfill
  Check your settings:
\begin{lstlisting}
git config --list
\end{lstlisting}
\end{frame}

% ------------------------------------------------------------------
\section{Creating a repository}

\begin{frame}[fragile]{Starting a repository}
  A \emph{repository} (``repo'') is a folder whose history Git tracks.
\begin{lstlisting}
mkdir my-paper
cd my-paper
git init
\end{lstlisting}
\vfill
This creates a hidden folder \texttt{.git/} where Git stores everything.

% \vfill
  \alert{Do not touch \texttt{.git/} by hand.}
\end{frame}

\exercise{
  Create a new repository
}


% ------------------------------------------------------------------
\section{Existing repositories}

\begin{frame}[fragile]{Cloning}
  To get a copy of an existing repository (notice the \texttt{.git} at
  the end):
  \begin{lstlisting}
    git clone https://github.com/user/repo.git
  \end{lstlisting}
  This:
  \begin{itemize}
    \item downloads the full history
    \item creates a folder \texttt{repo/}
    \item sets up a connection to the original (the \emph{remote})
  \end{itemize}
  \vfill
  You now have your own, fully independent copy.
\end{frame}

\exercise{
  Clone the repository of this course (note: this is not
  repository url)

  \url{https://github.com/davidemanini/git-course-2026}
}

\section{Daily workflow}

\begin{frame}{The three areas}
  Git works with three places:
  \vfill
  \begin{center}
  \begin{tabular}{ll}
    \textbf{Working directory} & the files you actually edit \\
    \textbf{Staging area} (index) & what will go in the next commit \\
    \textbf{Repository} & the history of committed snapshots \\
  \end{tabular}
  \end{center}
  \vfill
  The basic cycle:
  \[
    \text{edit} \longrightarrow \texttt{git add <files>} \longrightarrow \texttt{git commit}
  \]
\end{frame}

\begin{frame}[fragile]{The basic cycle}
\begin{lstlisting}
git status                    # what is going on?
# edit paper.tex
git add paper.tex             # stage a file
git add .                     # stage everything
git commit                    # record a snapshot
\end{lstlisting}
  \vfill
  \texttt{git status} is your best friend. Run it often.
\end{frame}

\begin{frame}{Writing good commit messages}
  A commit message should answer: \emph{what changed and why?}
  \vfill
  Conventions:
  \begin{itemize}
    \item Short first line (${\leq}$ 70 characters)
    \item Use the imperative: ``Add lemma\ldots'', not ``Added\ldots''
%    \item Optional longer body after a blank line
  \end{itemize}
  \vfill
  {\color{blue}Good} and {\color{red} bad} examples:
  \begin{itemize}
  \item \texttt{Fix proof of Theorem
      {\color{red}2.1}
    {\color{blue} on triangular comparison}}
    \item \texttt{Add section on Sobolev embeddings}
    \item \texttt{Rewrite introduction after referee report}
      \item \texttt{\color{red}Add details}
  \end{itemize}
\end{frame}

\begin{frame}
  Key takeaways: \LARGE 
  \begin{itemize}
  \item Commits are {\tiny almost always} immutable:

    think before committing
  \item Good commits messages are paramount
  \item Commits should be atomic
  \end{itemize}
\end{frame}


\exercise{
  \begin{itemize}
  \item
    Create a few new commits
  \end{itemize}
}
% ------------------------------------------------------------------
\section{Inspecting history}

\begin{frame}[fragile]{Looking at the history}
\begin{lstlisting}
git log
git log --oneline
\end{lstlisting}
%git show <commit-hash>
  \vfill
  Every commit has a unique hash (\texttt{a3f1c9e\ldots}). The first
  7 characters are usually enough.
\end{frame}

\begin{frame}[fragile]{Seeing what changed}
\begin{lstlisting}
git diff                 # unstaged changes
git diff --staged        # staged but not committed
git diff <commit>        # vs. a past commit
git diff <old> <new>     # compare two commits
\end{lstlisting}
  \vfill
  Output uses the standard ``unified diff'' format:
  lines starting with \texttt{-} are removed, \texttt{+} are added.
\end{frame}

\exercise{
  \begin{itemize}
  \item Play with \texttt{git log} and \texttt{git diff}
  \end{itemize}
}

% ------------------------------------------------------------------
\section{Ignoring files}

\begin{frame}[fragile]{\texttt{.gitignore}}
  Not every file should be tracked: compiled output, backups, caches\ldots
  \vfill
  Create a file called \texttt{.gitignore} in the repo root:
\begin{lstlisting}
# LaTeX auxiliary files
*.aux
*.log
*.out
*.synctex.gz
*.toc
*.pdf
\end{lstlisting}
\vfill
The \texttt{.gitignore} of this repository is appropriate for \LaTeX\
projects

\vfill
Takeaway: adding \texttt{.gitignore} should be one of the first commits
\end{frame}

\exercise{
  Add a file \texttt{.gitignore} to your project and commit it
}


\begin{frame}{Summary}
  You now know how to:
  \begin{itemize}
    \item Move around a shell
    \item Create a Git repository
    \item Stage and commit changes
    \item Inspect history and diffs
    \item Clone existing repositories
    \item Ignore unwanted files
  \end{itemize}
  \vfill
  Next lecture: branches, merging, and GitHub.
\end{frame}

% ------------------------------------------------------------------

% \begin{frame}{Exercise}
%   \begin{enumerate}
%     \item Create a directory \texttt{workshop/} and \texttt{cd} into it.
%     \item Run \texttt{git init}.
%     \item Create a small \texttt{hello.tex} (or \texttt{hello.py}).
%     \item Commit it.
%     \item Edit it, run \texttt{git diff}, then commit the change.
%     \item Add a \texttt{.gitignore} and commit it.
%     \item Inspect the history with \texttt{git log -{}-oneline}.
%   \end{enumerate}
%   \vfill
% %  Try to break things and ask questions.
% \end{frame}

\section{Hausaufagabe}

\exercise{
  \begin{itemize}
  \item Familiarize with the shell
  \item Setup git to a different computer
  \item Create a new repository a make some commits
  \item Inspect history and compare commits
  \end{itemize}
}

\end{document}
